Three lines through the corners of any triangle, trisecting its angles, meet in the corners of an
equilateral triangle, and nobody noticed until 1899. Circles packed into a circle can have whole-number
curvatures forever. A chain of circles between two others closes up or never does, wherever you start. Six
points on an ellipse always put three intersection points on one line. Geometry is full of such miracles, and
this book explains every one of them.

It starts in the plane of Euclid, seen through complex numbers, and travels through symmetry groups,
inversion, the projective plane, conics, quadrics and plane curves to the cubic curves whose points can be
added, the arithmetic heart of modern number theory. Then come lattices and the densest packings of balls,
the curved worlds of the sphere and the hyperbolic plane, the inequalities of convex bodies, and curvature,
from the four-vertex theorem to Gauss--Bonnet. Every theorem is proved, every solution explains how to find
the idea, and a Python notebook brings each chapter to life.
